% 364_Ikosaeder_D4_Vergleich_En_ch.tex
% Johann Pascher, 14 Sept. 2026
% The icosahedral carrier and the D4 lattice: Galois structure, reduction mod 3 and the doubling 3+3'

\providecommand{\BM}{\textbf{[B]}}
\providecommand{\KM}{\textbf{[K]}}
\providecommand{\SM}{\textbf{[S]}}
\providecommand{\SetzM}{\textbf{[POSIT]}}
\providecommand{\QM}{\textbf{[Q]}}
\providecommand{\EM}{\textbf{[E]}}
\providecommand{\XM}{\textbf{[X]}}

\chapter*{The Icosahedral Carrier and the D\textsubscript{4} Lattice}
\addcontentsline{toc}{chapter}{The Icosahedral Carrier and the D₄ Lattice}

\begin{abstract}
An externally proposed carrier model for emergent physics uses a
twelve-port icosahedron with symmetry group $A_5$ and a local averaging
repair rule on its thirty edges. This document compares that carrier
exactly with the FFGFT geometry $T^4/\mathbb{Z}_3$ on the $D_4$ lattice.
Four verification scripts establish: (i)~the carrier spectrum lies in
$\mathbb{Q}(\sqrt5)$, and the Galois conjugation $\sqrt5\mapsto-\sqrt5$
exchanges the two three-dimensional $A_5$ blocks $3\leftrightarrow3'$ as
spectral projectors, so that over $\mathbb{Q}$ only the 6D block exists --
the doubling $6=3\oplus3'$ of Dok.~285; (ii)~the reduction
$\mathbb{Z}[\varphi]/(3)=GF(9)$ turns this Galois conjugation into the
Frobenius $x\mapsto x^3$ of the FFGFT ground field; (iii)~$A_5$ acts
faithfully and absolutely irreducibly on $GF(9)^3$ but not on $GF(9)^2$,
hence not in $G(3)$ but in $G(6)$ -- the same threshold as the lepton
order~26 in Dok.~341; (iv)~$A_5$ is incompatible with the $D_4$ lattice
($5\nmid1152$, Dok.~314); the common core is the tetrahedral group $A_4$,
which acts simply transitively on the twelve ports, and the port
$\mathbb{Z}_3$ is provably never one of the sixteen orbifold generators of
$T^4/\mathbb{Z}_3$. The gauge assignment of the external model
($u(1)\oplus su(2)\oplus su(3)$ on $1\oplus3\oplus(3'\oplus5)$) is
representation-theoretically consistent; its Lie structure is an axiom
there, not a consequence of the icosahedron.
\end{abstract}

\section*{List of symbols}
\addcontentsline{toc}{section}{List of symbols}

The table explains symbols appearing in more than one section.
Quantities used only locally are introduced there.

\medskip
\begin{tabular}{@{}lp{11cm}@{}}
\toprule
\textbf{Symbol} & \textbf{Meaning} \\
\midrule
\multicolumn{2}{l}{\textit{Graph and carrier}} \\
$n$ & number of ports; $n=12$ (icosahedron) \\
$A$ & adjacency matrix; $A_{ij}=1$ if seam $(i,j)$ \\
$L=5I-A$ & Laplacian; 5-regular graph \\
$W=A/5$ & averaging operator (expected single-seam update) \\
$\mu$ & Laplacian eigenvalue; $\mu\in\{0,\,5{-}\sqrt5,\,6,\,5{+}\sqrt5\}$ \\
\midrule
\multicolumn{2}{l}{\textit{Groups and representations}} \\
$A_5$ & icosahedral rotation group; $|A_5|=60$ \\
$A,T_1,T_2,H$ & irreducible $A_5$ representations; dimensions $1,3,3,5$ \\
$\varphi=(1{+}\sqrt5)/2$ & golden ratio \\
$A_4$ & tetrahedral group; integral elements of $A_5$; $|A_4|=12$ \\
${\rm Aut}(D_4)$ & automorphism group of the $D_4$ lattice; $|{\rm Aut}(D_4)|=1152$ \\
$W(F_4)$ & Weyl group of $F_4$; isomorphic to ${\rm Aut}(D_4)$ \\
\midrule
\multicolumn{2}{l}{\textit{Galois and fields}} \\
$\mathbb{Q}(\sqrt5)$ & number field; generated by the carrier spectrum \\
$GF(3^k)$ & finite field of characteristic~3 \\
$\sigma:\sqrt5\mapsto-\sqrt5$ & Galois automorphism of $\mathbb{Q}(\sqrt5)$ \\
$x\mapsto x^3$ & Frobenius automorphism of $GF(9)$ \\
$\zeta_n$ & primitive $n$-th root of unity \\
\midrule
\multicolumn{2}{l}{\textit{$D_4$ lattice}} \\
$D_4$ & root lattice; $\{k\in\mathbb{Z}^4:\sum k_i\equiv0\pmod2\}$; kissing number~24 \\
$P_{T_1},P_{T_2},P_A,P_H$ & spectral projectors of $L$ \\
\bottomrule
\end{tabular}


% ============================================================
\section*{Status markers}
% ============================================================
\begin{center}
\begin{tabular}{@{}lp{10cm}@{}}
\textbf{[B]} & Algebraically proven (verification script, exact arithmetic). \\[3pt]
\textbf{[K]} & Numerically confirmed. \\[3pt]
\textbf{[S]} & Structurally plausible, physical identification open. \\[3pt]
\textbf{[E]} & Established external mathematics. \\[3pt]
\textbf{[X]} & Checked and negative. \\
\end{tabular}
\end{center}

This document contains no new mathematics. The character table of $A_5$,
the Galois theory of $\mathbb{Q}(\sqrt5)$, Brauer degrees, ${\rm Aut}(D_4)$,
the crystallographic restriction and cut-and-project are textbook material
\EM. The marker \BM\ here denotes the exact verification of such results
on the concrete carrier and their assembly; only the application to the
relation between the two frameworks is new.

% ============================================================
\section{Orientation: geometry, algebra, and what a field is here}
% ============================================================

FFGFT has two layers that must not be confused.

\paragraph{The geometry} is a torus. The carrier is $T^4$ with the $D_4$
lattice (kissing number~24, Dok.~314), and the physics sits in the
\emph{modes} of this torus: the winding numbers of the plane waves living
on it. The built-in $\mathbb{Z}_3$ symmetry of the lattice turns it into
the orbifold $T^4/\mathbb{Z}_3$ with nine fixed points. What a particle is,
which mass it carries, which sector it belongs to -- these are statements
about modes and their $\mathbb{Z}_3$ orbits.

\paragraph{The algebra} consists of finite fields: $GF(9)$ as ground
field, $GF(27)$ with its 26 units in eight three-orbits, both sitting in
parallel inside $GF(729)$ (Dok.~338, 341). The Frobenius $x\mapsto x^3$ is
the Galois group of these fields, and the sector pairing
$k\leftrightarrow-k$ is its action.

\paragraph{How the layers connect.} The fields \emph{are} not the modes.
They are the eigenvalue fields, splitting fields and coefficient fields in
which the mode operators live, and they decide which modes \emph{can
exist}: a mode is an element order in a matrix algebra over $GF(3^k)$, and
Lagrange permits only orders dividing the group order $|GF(3^k)^*|$. That
is why order~5 exists in $G(3)$ ($5\mid80$) while order~26 appears only in
$G(6)$ ($13\nmid80$, Dok.~341). In short: \emph{the geometry supplies the
modes; the Galois structure of the fields says which of them are
admissible.} The fields are mode-conditioned in exactly this sense -- they
are called up by the modes, not the other way round.

\paragraph{Why this organises the comparison.} OPH has the same two-layer
structure with different entries: the geometry is a sphere of icosahedra,
the modes are the blocks of the port space under $A_5$, and the algebra is
the number field $\mathbb{Q}(\sqrt5)$ with its Galois group
$\sqrt5\mapsto-\sqrt5$. There too the Galois structure decides which modes
exist over the ground field: over $\mathbb{Q}$ only the 6D block, never a
single triplet. The comparison in this document is therefore a comparison
of the \emph{algebras that classify the modes} -- and the two algebras
touch exactly once: in the reduction modulo~3, which maps
$\mathbb{Q}(\sqrt5)$ onto $GF(9)$ and the Galois conjugation onto the
Frobenius (Theorem~\ref{thm:frob}). The geometries, by contrast, touch only
in the tetrahedral group and part at the five-fold symmetry
(Section~\ref{sec:d4}).

% ============================================================
\section{Occasion and scope}
% ============================================================

In September 2026 a carrier model was presented on the IPI list (Observer
Patch Holography, hereafter \emph{OPH}) whose hardware is a finite network
of twelve-port icosahedra and whose software is a single local rule: on
each of the thirty edges (\emph{seams}) the two port readings are replaced
by their arithmetic mean. The author supplied a reproducibility bundle in
exact rational arithmetic (300 updates, seeds 7 and 11, an independent
verifier).

This document takes no position on the physical claims of OPH. It answers
three narrow, exactly checkable questions:
\begin{enumerate}[nosep]
\item What is the structure of the carrier spectrum, and what does its
      Galois group imply for the distinguishability of the two
      three-dimensional blocks?
\item How does $\mathbb{Q}(\sqrt5)$ relate to the FFGFT ground field $GF(9)$?
\item Which part of the icosahedral group $A_5$ is compatible with the
      $D_4$ lattice and the orbifold $T^4/\mathbb{Z}_3$?
\end{enumerate}
The FFGFT reference points are the doubling $6=3\oplus3'$ (Dok.~285), the
embedding $C_3<A_5$ (Dok.~293), the negative result $5\nmid|{\rm Aut}(D_4)|$
(Dok.~314) and the $G(3)/G(6)$ threshold (Dok.~341).

% ============================================================
\section{The spectrum of the icosahedral carrier}
% ============================================================

\begin{theorem}[Spectrum, \EM/\BM]\label{thm:spektrum}
Let $A$ be the adjacency matrix of the icosahedral graph (12 vertices,
30 edges, 5-regular) and $L=5I-A$ its Laplacian. Then
\[
  \operatorname{spec}(A)=\{5^{(1)},\ \sqrt5^{(3)},\ -1^{(5)},\ -\sqrt5^{(3)}\},\qquad
  \operatorname{spec}(L)=\{0^{(1)},\ (5-\sqrt5)^{(3)},\ 6^{(5)},\ (5+\sqrt5)^{(3)}\}.
\]
The only rational non-trivial Laplacian eigenvalue is $6$ (multiplicity~5).
The field generated by the spectrum is $\mathbb{Q}(\sqrt5)$, of degree~2.
\end{theorem}

\begin{proof}
Script \texttt{spectrum\_icosahedral\_carrier.py}, SymPy, exact; the seam
list is taken from the OPH bundle, alternatively from the canonical
vertices $(0,\pm1,\pm\varphi)$ cyclically. The script also replays the 300
updates of seed~7 in rational arithmetic and reproduces the centered
residual recorded in the bundle exactly.
\end{proof}

The multiplicities $(1,3,5,3)$ are the dimensions of the irreducible
$A_5$ representations $A\oplus T_1\oplus H\oplus T_2$ (icosahedral-group
notation; $H$ is five-dimensional, $T_1,T_2$ are the two inequivalent
triplets). On a five-fold rotation $g_5$ they carry the character values
\[
  \chi_{T_1}(g_5)=\varphi,\qquad \chi_{T_2}(g_5)=1-\varphi,\qquad
  \chi_{T_1}+\chi_{T_2}=1 .
\]

\begin{remark}[Damping rates, \EM/\BM]
The expected action of a random single-seam update is $E[M]=I-L/60$ with
block eigenvalues $(55+\sqrt5)/60\approx0.954$ on $T_1$, $9/10$ on $H$ and
$(55-\sqrt5)/60\approx0.879$ on $T_2$. $T_1$ is the slowest block; its
distance to $T_2$ is $\sqrt5/30$ (ratio $\approx1.085$). This is a
statement about rates, not a symmetry selection.
\end{remark}

% ============================================================
\section[Galois conjugation and the doubling 6 = 3 + 3']{Galois conjugation and the doubling $6=3\oplus3'$}
% ============================================================

\begin{theorem}[Galois exchanges $T_1\leftrightarrow T_2$, \EM/\BM]\label{thm:galois}
Let $P_{T_1},P_{T_2},P_A,P_H$ be the spectral projectors of $L$
(polynomials in $L$ with coefficients in $\mathbb{Q}(\sqrt5)$). The Galois
automorphism $\sigma:\sqrt5\mapsto-\sqrt5$ of $\mathbb{Q}(\sqrt5)$ acts
entrywise as
\[
  \sigma(P_{T_1})=P_{T_2},\qquad \sigma(P_{T_2})=P_{T_1},\qquad
  \sigma(P_A)=P_A,\qquad \sigma(P_H)=P_H .
\]
The projector $P_{T_1}+P_{T_2}$ is $\mathbb{Q}$-rational; $P_{T_1}$ and
$P_{T_2}$ separately are not. The minimal polynomial of the $T_1$
eigenvalue over $\mathbb{Q}$ is $x^2-10x+20$ with roots $5\mp\sqrt5$.
\end{theorem}

\begin{proof}
Script \texttt{galois\_z3\_icosahedral\_carrier.py}. The projectors are
built as Lagrange polynomials in $L$; the exchange is checked by
substituting $\sqrt5\to-\sqrt5$ and comparing matrices.
\end{proof}

\begin{corollary}[Indivisibility over $\mathbb{Q}$, \BM]\label{cor:sechs}
Every $A_5$-invariant structure on the port space that is defined over
$\mathbb{Q}$ \emph{as an abstract field} -- a lattice, an integer
character -- sees $T_1\oplus T_2$ as \emph{one} six-dimensional block;
$T_1\oplus T_2$ is the only six-dimensional submodule with integer
character that contains a triplet ($1\oplus H$ is the other six-dimensional
submodule with integer character). Distinguishing $T_1$ from $T_2$ is not a
symmetry statement.
\end{corollary}

\begin{remark}[The selection happens at the archimedean place, \BM]\label{rem:arch}
Argument of the author of the external model (personal communication,
15~September 2026), verified exactly here.
The repair dynamics does not live in the abstract field $\mathbb{Q}$ but in
$\mathbb{R}$: $I-L/60$ is a rational symmetric matrix with real
eigenvalues, and the two conjugate rates $(55\pm\sqrt5)/60$ are the real
numbers $0.954$ and $0.879$. Both embeddings
$\mathbb{Q}(\sqrt5)\hookrightarrow\mathbb{R}$ yield the same two subspaces;
they differ only in which one carries the name $T_1$. Iterating $I-L/60$ on
any rational vector with a component in the slow block, the ratios of the
centered norms -- a rational sequence -- tend to
$((55+\sqrt5)/60)^2=0.9100$ and not to the conjugate value $0.7733$
(script, section~6). Selecting the slow triplet as space is therefore a
derivation from the dynamics, as every convergent rational sequence
selects the positive root of $x^2=2$; it happens at the archimedean place
of $\mathbb{Q}(\sqrt5)$. Geometrically: in the coordinates
$(0,\pm1,\pm\varphi)$ the five seam neighbours of a port sum to
$\sqrt5\,v$, the thirty second neighbours (edges of the great icosahedron)
to $-\sqrt5\,v$; \enquote{$T_1$ is space} says that repair couples nearest
neighbours.
\end{remark}

This is the theorem behind Dok.~285: there the doubling $6=3\oplus3'$ is
the answer to the crystallographic restriction (no icosahedral lattice in
three dimensions, but one in six). Here the same statement appears as
Galois indivisibility at the \emph{finite} places; the archimedean place
separates the triplets, the place~3 does not (Theorem~\ref{thm:frob},
Section~\ref{sec:phasen}). OPH's proposal that the three-dimensional continuum
is the dense projection of \enquote{integer frames built from port labels}
is the quasicrystal projection $\mathbb{Z}^6\to\mathbb{R}^3$; Dok.~285
describes the cause, OPH the image.

\begin{theorem}[Restriction to $C_3<A_5$, \EM/\BM]\label{thm:c3}
Under a three-fold rotation $g_3\in A_5$ the port space splits into
$\mathbb{Z}_3$ eigenspaces of dimensions $4+4+4$; blockwise
\[
  A\to 1,\qquad T_1\to 1\oplus\omega\oplus\omega^2,\qquad
  T_2\to 1\oplus\omega\oplus\omega^2,\qquad H\to 1\oplus2\omega\oplus2\omega^2 .
\]
The trivial part has dimension~4, the non-trivial part dimension~8.
\end{theorem}

This is the $C_3$-in-$A_5$ embedding of Dok.~293 in character form. The
non-trivial part (dimension~8) is, as an abstract $\mathbb{Z}_3$-module,
two copies of the orbifold representation on $\mathbb{R}^4$; whether more
than representation theory is contained in this is settled negatively in
Section~\ref{sec:d4}.

% ============================================================
\section[{Reduction modulo 3: Z[phi]/(3) = GF(9)}]{Reduction modulo 3: $\mathbb{Z}[\varphi]/(3)=GF(9)$}
% ============================================================

\begin{theorem}[Inert prime, \EM/\BM]\label{thm:inert}
$x^2-x-1$ is irreducible over $GF(3)$ ($5\equiv2$ is not a square
modulo~3). Hence $3$ is inert in $\mathbb{Q}(\sqrt5)$ and
$\mathbb{Z}[\varphi]/(3)\cong GF(9)$.
\end{theorem}

\begin{theorem}[Galois becomes Frobenius, \EM/\BM]\label{thm:frob}
Under the reduction of Theorem~\ref{thm:inert} the Galois conjugation
$\varphi\mapsto1-\varphi$ becomes the Frobenius $x\mapsto x^3$ of $GF(9)$:
in $GF(9)$ one has exactly $\varphi^3=1-\varphi$.
\end{theorem}

\begin{proof}
$\varphi^2=\varphi+1$, hence $\varphi^3=2\varphi+1\equiv1-\varphi\pmod3$.
Script \texttt{zphi\_mod3\_a5.py}, sections (1)--(2).
\end{proof}

Thus the FFGFT sector pairing (Frobenius on $GF(27)^*$, Dok.~343/363) and
the exchange $T_1\leftrightarrow T_2$ of the icosahedral carrier are not
merely analogous: modulo~3 they are the same map.

\begin{theorem}[$A_5$ over $GF(9)$, \EM/\BM]\label{thm:a5gf9}
The sixty rotation matrices of the icosahedron (entries in
$\tfrac12\mathbb{Z}[\varphi]$) reduce modulo~3 to sixty distinct matrices in
$GL_3(GF(9))$ spanning all of $M_3(GF(9))$: the reduction is faithful and
absolutely irreducible. Entrywise Frobenius of the reduced $T_1$ matrices
yields exactly the reduced $T_2$ matrices. By contrast $GL_2(GF(9))$
(order 5760, searched exhaustively) contains no $(2,3,5)$-generated
subgroup of order~60: $A_5$ does not embed in $GL_2(GF(9))$.
\end{theorem}

\begin{corollary}[The $G(3)/G(6)$ threshold, \SM]
With $G(3)\cong M_2(GF(9))^2$ and $G(6)\cong M_8(GF(9))$ (Dok.~341), $A_5$
is not in the unit group of $G(3)$ but is in that of $G(6)$. This is the
same threshold at which Dok.~341 finds the lepton order~26. The reasons
are formally different (Brauer degrees $\{1,3,3,4\}$ of $A_5$ in
characteristic~3 versus Lagrange in $GF(81)^*$); whether the coincidence
has substance remains open.
\end{corollary}

% ============================================================
\section[A5 versus Aut(D4)]{$A_5$ versus ${\rm Aut}(D_4)$}\label{sec:d4}
% ============================================================

\begin{theorem}[Reproduction of Dok.~314, \EM/\BM]\label{thm:aut}
The automorphism group of the $D_4$ root system, recomputed from the 24
roots, has order $1152$ with order distribution
$(1,2,3,4,6,8,12)\mapsto(1,139,80,228,464,144,96)$ and no element of
order~5. Among the 80 elements of order~3 exactly 16 are fixed-point-free
on $\mathbb{R}^4$ ($|\det(1-A)|=9$); each acts freely on the 24 roots
(eight orbits of three).
\end{theorem}

\begin{corollary}[\XM]
$A_5$ acts on no $D_4$ lattice: $5\nmid1152$. The symmetry group of OPH
and the FFGFT lattice structure are incompatible as wholes.
\end{corollary}

\begin{theorem}[The common core is $A_4$, \EM/\BM]\label{thm:a4}
The elements of $A_5$ whose matrices in the model $(0,\pm1,\pm\varphi)$
are integral (no $\varphi$) form the tetrahedral group $A_4$ (orders
$1,2,3$ with counts $1,3,8$). $A_4$ acts simply transitively on the twelve
ports -- the port set is an $A_4$-torsor -- and embeds as
$\operatorname{diag}(M,1)$ in $W(D_4)<{\rm Aut}(D_4)$.
\end{theorem}

This sharpens Dok.~285 (\enquote{$C_3<A_5$}): not only the three-fold
symmetry but the whole tetrahedral group of order~12 -- the port count --
is lattice-compatible; only the five-fold symmetry drops out.

\begin{theorem}[Ports as half the first shell, \EM/\BM]
Every three-fold rotation of $A_5$ acts freely on the twelve ports (four
orbits of three). As a $\mathbb{Z}_3$-module the port space is
$4\cdot\mathbb{Z}[\mathbb{Z}_3]$, the first $D_4$ shell is
$8\cdot\mathbb{Z}[\mathbb{Z}_3]$: the ports are one half of
\enquote{$24=12+12$} from Dok.~314.
\end{theorem}

\begin{theorem}[The port $\mathbb{Z}_3$ is not an orbifold $\mathbb{Z}_3$, \XM]\label{thm:neg}
The three-fold rotation permuting the ports is a rotation of $\mathbb{R}^3$
and fixes an axis; as $\operatorname{diag}(M,1)$ it fixes a line in
$\mathbb{R}^4$ and is therefore one of the $80-16=64$ non-orbifold elements
of order~3. This is not an artefact of the embedding: for each of the 16
orbifold generators and each of the 70 choices of four of its eight root
orbits (1120 $\mathbb{Z}_3$-stable twelve-element halves) it was checked
whether the twelve roots form a cuboctahedron ($D_3$ half) or carry the
icosahedral graph on some inner-product class. Neither ever occurs.
\end{theorem}

\begin{proof}
A $\mathbb{Z}_3$ stabilising a $D_3$ half stabilises its normal up to sign,
hence exactly for order~3; it has a fixed vector and is not an orbifold
generator. Script \texttt{a5\_vs\_d4.py}, sections (E)--(F).
\end{proof}

\begin{remark}[Side result, \KM]
256 of the 1120 orbifold-stable twelve-element halves are tight frames in
$\mathbb{R}^4$ (Gram eigenvalue~6 with multiplicity~4, the 4D counterpart
of the OPH relation $G^2=4G$ for the port vectors). Their graph spectra
contain $\sqrt{13}$, not $\sqrt5$; the stabiliser has order~18 and is not
transitive. They carry no port structure.
\end{remark}

% ============================================================
\section{The OPH gauge assignment}
% ============================================================

OPH assigns to the port module the Lie algebra $u(1)\oplus su(2)\oplus su(3)$:
$u(1)$ on $A$, $su(2)$ on $T_1$, $su(3)$ on $T_2\oplus H$.

\begin{theorem}[Consistency of the assignment, \EM]
$A_5$ sits in $SO(3)\subset SU(3)$ via $T_2$. The adjoint of $su(3)$ is
$3\otimes\bar3-1$; under $SO(3)$ it splits as $3\oplus5$, under $A_5$ with
the $T_2$ embedding as $T_2\oplus H$ ($\Lambda^2T_2=T_2$,
$\operatorname{Sym}^2T_2=1\oplus H$). Likewise the adjoint of $su(2)$
under the $T_1$ embedding is $T_1$. The restriction of
$u(1)\oplus su(2)\oplus su(3)$ to $A_5$ is therefore exactly
$A\oplus T_1\oplus T_2\oplus H$.
\end{theorem}

\begin{remark}[What is axiom, \SM/\EM]
The derivation of the centre in this remark is due to the author of the
external model (personal communication, 15~September 2026).
The assignment presupposes two axioms of the external model: (A1) the
reversible port response is a compact Lie algebra $\mathfrak{g}$ with a
faithful representation on the port module; (A2) $A_5$ acts on
$\mathfrak{g}$ by inner automorphisms. The one-dimensional centre then
follows and is not a third axiom: inner automorphisms act trivially on the
centre, so it lies in the $A_5$-fixed part of $1\oplus3\oplus3'\oplus5$,
the uniform line. Centre zero would require $\dim\mathfrak{g}=12$
semisimple, i.e.\ $su(2)^4$; under an inner action each $su(2)$ factor has
fixed dimension $0$ or $3$, so the fixed dimension would be a multiple
of~3, whereas the port module has fixed dimension~1. Hence the centre is
the uniform line and $11=8+3$ is the only decomposition of the rest into
dimensions of compact simple algebras. The counterexample
$u(1)^3\oplus su(2)^3$ fails on (A2): under an inner action its module
contains no~$5$; and even without (A2) it fails by
Theorem~\ref{thm:galois}, since the centre of a torus is defined over
$\mathbb{Q}$ and can therefore contain $3$ and $3'$ only with equal
multiplicity. The correct statement is thus \enquote{$A_5$ plus (A1), (A2)
forces the Lie type}.
\end{remark}

FFGFT is in the mirror situation: $su(3)$ follows from the $\mathbb{Z}_3$
triality on $T^4$ \BM\ (Dok.~321), $su(2)$ is open \SM\ (Dok.~324). In OPH
$T_1$ is clean and $8=T_2\oplus H$ is the axiomatic step.

% ============================================================
\section{Phases as roots of unity: how the connection succeeds}\label{sec:phasen}
% ============================================================

Section~\ref{sec:d4} and Theorem~\ref{thm:frob} seem to contradict each other:
five-fold symmetry is incompatible with the lattice, yet after reduction
$A_5$ sits in $GL_3(GF(9))$. The contradiction dissolves once one
distinguishes \emph{as what} a symmetry enters the carrier -- as a lattice
automorphism or as a phase.

\begin{theorem}[Phase channel, \EM/\BM]\label{thm:phasen}
\begin{enumerate}[nosep]
\item Modulo~3, $\varphi$ is a primitive element of $GF(9)^*$:
      $\varphi^4=-1$, order~8.
\item The $72^\circ$ rotation has characteristic polynomial
      $(x-1)(x^2-(\varphi-1)x+1)$. Modulo~3 the quadratic factor has no root
      in $GF(9)$; its roots lie in $GF(81)$ and have order~5 there ($5\mid80$).
\item $13\mid728=|GF(729)^*|$ but $13\nmid80$: order~26 needs the cubic
      extension (Dok.~341).
\item In characteristic~3, $x^3-1=(x-1)^3$: no $GF(3^k)$ contains a
      primitive cube root of unity. A $\mathbb{Z}_3$ acts unipotently in the
      field, never as a phase~$\omega$.
\end{enumerate}
\end{theorem}

\begin{proof}
Script \texttt{phasen\_einheitswurzeln.py}.
\end{proof}

\paragraph{Reading.} A phase is a root of unity, and roots of unity are
exactly what a field extension adds. In the lattice the $72^\circ$
rotation is an automorphism of vectors, and $5\nmid1152$. In the field the
same rotation is only its eigenvalue $\zeta_5$, and $5\mid80$. Geometric
incompatibility becomes arithmetic compatibility as soon as one passes
from lattice vectors to phases. This is the mechanism behind Dok.~293:
$\theta=2/9$ arises from $R_5$ as a phase, not as a lattice operation.

The extension degrees are the two basic operations of FFGFT:
\begin{center}
\begin{tabular}{@{}llll@{}}
\toprule
Extension & Degree & carries & Counterpart \\
\midrule
$GF(9)\to GF(81)$ & 2 & $\zeta_5$, $\varphi$, $2/\sqrt5$, icosahedron & doubling $6=3\oplus3'$ (Dok.~285) \\
$GF(9)\to GF(729)$ & 3 & $\zeta_{13}$, order~26, leptons & $\mathbb{Z}_3$ triality (Dok.~341) \\
\bottomrule
\end{tabular}
\end{center}

And three-fold symmetry itself is \emph{not} a phase: in characteristic~3
$\omega$ collapses to $1$ (Dok.~341: the $N$ eigenvalues $\{-3,+3\}$ become
$\{0\}$). This is why in FFGFT the $\mathbb{Z}_3$ is lattice structure --
a circulant on the mode orbits, Dok.~314 -- and five-fold symmetry is a
phase. It is not a choice but forced: what cannot exist as a root of unity
in characteristic~3 must enter as a lattice operation; what does not fit
the lattice must enter as a root of unity in an extension. The two
channels are complementary, and the prime~3 decides which symmetry falls
into which channel. In the language of places of $\mathbb{Q}(\sqrt5)$: the
external model chooses between $T_1$ and $T_2$ at the archimedean place
(Remark~\ref{rem:arch}); FFGFT computes at the place~3, where $\sqrt5$ has
no sign and the two triplets are exchanged by the Frobenius. Both
selections are correct, at different places -- and that is why the two
frameworks touch exactly once. Theorem~\ref{thm:neg} is the geometric side of the
same statement: the port $\mathbb{Z}_3$ is lattice, the orbifold
$\mathbb{Z}_3$ is lattice, but they are different lattice operations; the
five-fold symmetry that distinguishes them is not present in the lattice
at all.

% ============================================================
\section{What is projected and what is not: scope of the conclusions}\label{sec:reichweite}
% ============================================================

The reference object of this document is the torus, not the icosahedron.
The icosahedron is \emph{one} geometric body among many that can be
projected onto $T^4/\mathbb{Z}_3$, and it is the test body here only
because an external model proposes it as carrier. Every other finite body
-- the five-fold geometry of Dok.~285 was the previous case -- splits by
the same pattern into a lattice-compatible part (in ${\rm Aut}(D_4)$) and
a phase part (roots of unity in an extension $GF(3^k)$); the icosahedron
is the instance on which the pattern is worked out completely in this
document, not the pattern itself. None of the maps is a complete
projection of the icosahedron onto the torus; by Theorem~\ref{thm:aut} no
such projection can exist, and that holds for every body with
non-crystallographic symmetry. In the terminology of Dok.~330 the maps are
neither Type~I (lossless decompactification) nor Type~III (bijective
representational translation), but lossy like Type~II -- or not FFGFT
operations at all. The conclusions are justified exactly as far as the
respective map is faithful on what is being concluded. They also hold
relative to the choice of lattice: $D_4$ is a posit (Dok.~320 \SetzM),
motivated by the densest packing in four dimensions and by the fact that
$D_4$ is the only lattice of the $D_n$ family with $S_3$ outer symmetry,
which carries the $\mathbb{Z}_3$ as triality (Dok.~314; there also the
Casimir counter-finding in favour of $\mathbb{Z}^4$). All statements of
this document about ${\rm Aut}(D_4)$ presuppose that choice.

\paragraph{Genuine projections (faithful on the object of the conclusion).}
\begin{enumerate}[nosep]
\item $\mathbb{Z}_3$ on the mode orbits (Dok.~314): intrinsic torus
      structure, complete. All triality conclusions rest on it.
\item $\mathbb{Z}[\varphi]\to GF(9)$ (Theorem~\ref{thm:inert}): surjective
      ring homomorphism with kernel $(3)$. Faithful for everything at the
      place~3 (mode orders, extension degrees, indistinguishability of
      $T_1,T_2$); lost is the sign of $\sqrt5$ and with it everything
      archimedean (rates, order, convergence). No conclusion about dynamics.
\item Rotation $\mapsto$ eigenvalue (Theorem~\ref{thm:phasen}): faithful for
      the order, with maximal kernel -- the geometry is forgotten. The
      conclusion \enquote{five-fold symmetry only as a phase} holds for
      precisely that reason.
\item $A_4\hookrightarrow W(D_4)$ (Theorem~\ref{thm:a4}): faithful for
      everything tetrahedral, five-fold symmetry not included.
\end{enumerate}

\paragraph{Not projections of the torus.}
\begin{enumerate}[nosep,resume]
\item Ports $=$ half the first shell: isomorphic only as $\mathbb{Z}_3$-sets;
      the icosahedral edges do not exist in $D_4$. Counting statements are
      justified, statements about neighbourhood or rates are not.
\item Cut-and-project $\mathbb{Z}^6\to\mathbb{R}^3$: a genuine projection of a
      \emph{different} lattice ($\mathbb{Z}^6/D_6$, the second protected
      geometry of Dok.~285), not of $T^4$.
\item Port $\mathbb{Z}_3\to$ orbifold $\mathbb{Z}_3$: provably not a
      projection (Theorem~\ref{thm:neg}).
\item Gauge assignment and choice of space in the external model:
      statements at the archimedean place about its response algebra;
      none of it projects onto $T^4$.
\end{enumerate}

\paragraph{What kind of confirmation this is.} The projections do not
confirm that the torus is the physics. They confirm that the torus
geometry makes concrete, falsifiable structural predictions and that an
independently constructed foreign object satisfies them. Four predictions
had the opportunity to fail:
\begin{itemize}[nosep]
\item \emph{Lattice}: order~5 does not act on $D_4$ (Dok.~314). Had
      $A_5\subset{\rm Aut}(D_4)$, Dok.~285 would be wrong. Passed.
\item \emph{Field}: $\zeta_5$ only in $GF(81)$, $\zeta_{13}$ only in
      $GF(729)$, no $\zeta_3$ (Dok.~341). Had $3$ split in $\mathbb{Q}(\sqrt5)$
      or $\zeta_5\in GF(9)$, the correspondence doubling $\leftrightarrow$
      degree~2 would be wrong. Passed.
\item \emph{Orbifold}: $\mathbb{Z}_3$ is a lattice operation, not a phase.
      The port $\mathbb{Z}_3$ could have been an orbifold $\mathbb{Z}_3$; it is
      not, for a reason that mirrors the torus structure. Passed.
\item \emph{Mode order}: the $G(3)/G(6)$ threshold meets the same boundary
      for $A_5$ as for order~26. The weakest test, \SM.
\end{itemize}
This is confirmation in the sense of a survived test: more than internal
consistency, less than proof -- and it applies to the torus, not to the
icosahedron. A further test body (octahedron, cube, a four-dimensional
polytope) would put the same four predictions to the question again, with
a different lattice part and a different number field; the tetrahedral and
octahedral groups lie entirely in ${\rm Aut}(D_4)$ and would have no phase
part, the $600$-cell would bring back $\mathbb{Q}(\sqrt5)$. In this series
the icosahedron is only the first fully computed case. Unexamined -- and not examinable by any
projection -- remains whether $T^4$ is the physical carrier, whether
$\theta=2/9$ is physically real as a phase (still \KM, Dok.~293) and
whether mode orders are particles; that requires measurement. Equally
untouched is the deeper question whether FFGFT's own translations
(Type~III, Dok.~270/330) are complete: there, internal consistency without
any statement about observability continues to hold. Whoever reads more
into this document than the location of the single point of contact
overreaches.

\section*{Tool and method}
\addcontentsline{toc}{section}{Tool and method}

Claude (Anthropic) was used as a tool for: writing and executing the Python
verification scripts; drafting the \LaTeX\ sources from results and corrections
of the author; maintaining the changelog and cross-references. The tool computes
and writes; the author decides what is computed and written, judges every result
and determines what it means physically. Every algebraic statement is backed by
a Python verification script with explicit assertions --- if an assertion fails,
the statement is not used.

% ============================================================
\section{Summary}
% ============================================================

\begin{center}
\begin{tabular}{@{}lllp{4.6cm}@{}}
\toprule
 & OPH & FFGFT & Link \\
\midrule
Field & $\mathbb{Q}(\sqrt5)$ & $GF(9)$, $GF(27)\parallel GF(729)$ & $\mathbb{Z}[\varphi]/(3)=GF(9)$; Galois $\to$ Frobenius \BM \\[2pt]
Group & $A_5$ & ${\rm Aut}(D_4)$, $\mathbb{Z}_3$ & $A_5\not\subset{\rm Aut}(D_4)$; core $A_4$ (= 12 ports) \BM \\[2pt]
$3\oplus3'$ & $T_1$ space, $T_2$ colour & 6D lattice (Dok.~285) & indivisible over $\mathbb{Q}$ \BM \\[2pt]
$\mathbb{Z}_3$ orbits & 4 free on 12 & 8 free on 24 & ports = half the first shell \BM \\[2pt]
Port $\mathbb{Z}_3$ & axis in $\mathbb{R}^3$ & orbifold fixed-point-free & never conjugate \XM \\[2pt]
$G(3)/G(6)$ & $A_5\not\subset GL_2(GF(9))$, $\subset GL_3$ & order~26 only in $G(6)$ & same threshold \SM \\
\bottomrule
\end{tabular}
\end{center}

\paragraph{Verification scripts.} \texttt{2/python/Dok364\_Skripte/}:
\texttt{spectrum\_icosahedral\_carrier.py},
\texttt{galois\_z3\_icosahedral\_carrier.py},
\texttt{zphi\_mod3\_a5.py}, \texttt{a5\_vs\_d4.py}, \texttt{phasen\_einheitswurzeln.py}; batch run
\texttt{run\_all\_364.sh}. The seam list of the external bundle is read if
present; otherwise the canonical icosahedron is used. No register entry:
no earlier document is corrected.

\begin{thebibliography}{9}

\bibitem{Conway1993}
J.~H.~Conway, N.~J.~A.~Sloane,
\textit{Sphere Packings, Lattices and Groups},
3rd ed., Springer, 1993.
Standard reference for Aut$(D_4)$, kissing number~24, order distribution.

\bibitem{Humphreys1990}
J.~E.~Humphreys,
\textit{Reflection Groups and Coxeter Groups},
Cambridge University Press, 1990.
$D_4$ triality, Dynkin diagram, $W(F_4)$.

\bibitem{Serre1977}
J.-P.~Serre,
\textit{Linear Representations of Finite Groups},
Springer, 1977.
Character table of $A_5$; Brauer degrees in characteristic~3.

\bibitem{Lidl1997}
R.~Lidl, H.~Niederreiter,
\textit{Finite Fields},
2nd ed., Cambridge University Press, 1997.
$GF(9)$ as extension of $GF(3)$; Frobenius automorphism; roots of unity in $GF(3^k)$.

\bibitem{Neukirch1999}
J.~Neukirch,
\textit{Algebraic Number Theory},
Springer, 1999.
Inertness of~$3$ in $\mathbb{Q}(\sqrt5)$; Galois group of $\mathbb{Q}(\sqrt5)/\mathbb{Q}$.

\bibitem{Berger1987}
M.~Berger,
\textit{Geometry~I},
Springer, 1987.
Icosahedron, $A_5$ symmetry; canonical vertices $(0,\pm1,\pm\varphi)$.

\bibitem{Pascher285}
J.~Pascher,
\textit{Dok.~285: FFGFT --- HLV, CP violation and dimension bridge},
FFGFT corpus.

\bibitem{Pascher293}
J.~Pascher,
\textit{Dok.~293: Icosahedron, $\theta=2/9$ and $C_3$-in-$A_5$},
FFGFT corpus.

\bibitem{Pascher314}
J.~Pascher,
\textit{Dok.~314: Lattice in Hilbert space},
FFGFT corpus.

\bibitem{Pascher341}
J.~Pascher,
\textit{Dok.~341: $GF(27)$ and GALG in FFGFT},
FFGFT corpus.

\end{thebibliography}
